% English abstract — set left-to-right inside the Persian document.

\chapter*{Abstract}
\addcontentsline{toc}{chapter}{Abstract (English)}
\markboth{Abstract}{Abstract}
\begin{latin}
\noindent
Diabetic retinopathy is the leading cause of preventable blindness among working-age adults, and
diabetic macular edema is the most common mechanism of central vision loss. Automated fundus
grading reports high headline figures, yet as of the end of April 2024 only three ophthalmic AI
devices had United States Food and Drug Administration clearance. This thesis begins from that
gap and asks what the reported numbers do and do not establish.

\vspace{0.6em}
\noindent\textbf{Methods.} A shared backbone with two threshold heads --- five ordinal grades for
retinopathy and three for macular edema --- was trained on 2{,}260 images from two public corpora
under five-fold grouped cross-validation. The split was frozen before the first real run and
built from perceptual clustering so that cross-corpus leakage is closed; 3{,}662 APTOS images were
held out until the final model was frozen. One rule governs every comparison in this work:
\textbf{no difference is attributed to representation unless both models are evaluated under
cross-fitted cut-points}. An untuned default is a hyper-parameter that has been chosen, not one
that has been avoided.

\vspace{0.6em}
\noindent\textbf{Results.} The selected model reaches a quadratic weighted kappa of 0.8954 on the
development pool and 0.9026 [0.8941, 0.9108] on the held-out external set, and the full pipeline
is worth +0.2342 kappa over frozen ImageNet features with a linear ordinal head on identical
images, folds, decoding and calibration. Of nine pre-registered interventions, eight were
falsified. The macular-edema head meets a ceiling that comes from data rather than architecture:
six independent interventions --- including a crop isolating the region the clinical grade is
defined on --- failed to move the confidence interval away from zero, and only 51 images in the
entire study carry the middle grade. A separate negative result establishes that the macular-edema
grade \textbf{cannot be recovered from published annotations}, closing the route to more
three-class data via segmentation masks.

\vspace{0.6em}
\noindent\textbf{Contribution.} The contribution is methodological rather than a performance
figure. The matched-calibration rule has reversed four prior attributions \emph{in both
directions}, with two sign reversals; more importantly, the same rule converted the model's
blindness to the mild grade from a capacity limit into a calibration problem solvable without
retraining. Alongside it, a set of data-provenance mechanisms is presented, motivated by five
worked examples --- \textbf{two of which are this thesis's own failures}: a chapter that
asserted its numbers were generated while they were typed by hand, three of them wrong and one
supporting a central argument with no artefact behind it at all; and, committed while writing
that up, a diagnosis taken from a specification rather than from the rendered page, which
``corrected'' something that was not broken.

\vspace{0.6em}
\noindent\textbf{Limitation.} The frozen, project-specific split makes every number here
incommensurable with any published figure. That trade is deliberate: comparability was given up
for correctness, and the reader should know which of the two they hold.

\vspace{1em}
\noindent\textbf{Keywords:} diabetic retinopathy, diabetic macular edema, ordinal grading,
calibration, decision threshold, external validation, data provenance, foundation models.
\end{latin}
\cleardoublepage
